\chapter{系统性能实验}
\label{ch:exp}

\section{实验目的与环境}

本章基于第二、三章给出的结构与参数设定，对已实现的哈希系统进行性能测试和功能验证。实验围绕操作延迟、k-kick 踢出次数、多层级 Cubby 向下拆分（rebuild\_down）与向上合并（rebuild\_up）三类指标展开，用于检查第三章各模块在实际运行中的效果\cite{bender2022otsh,bender2024sosa}。

实验的系统代码基于 C++23 标准开发实现，在 AMD Ryzen 7 6800H 处理器、Windows 10 环境下运行。元素规模 \(n\in\{10^3,10^4,10^5\}\)，Facility 规模 \(K\) 按分组取多档（基准 \(K=\log^3 n\)），k-kick 最大深度 \(k\in\{3,4,5\}\)（基准 \(k=4\)），NODE\_MAX\_BITS 常数 \(c=2\)。键由 std::mt19937\_64 在固定种子 seed=1（经 seed1/seed2/seed3 派生表内哈希参数）下均匀生成；实验报告单次 workload 内各操作类型的平均延迟。

\section{实验方案}

\subsection{实验分组与工作量}

实验分为四组，如表~\ref{tab:exp-plan} 所示。规模 \(n\) 组与 \(K\) 组、\(k\) 组采用键均匀分布于各 Facility 的混合增删负载；Cubby 动态重构组按纯插入、纯删除及不同插入/删除比例构造 workload，分别触发 rebuild\_up 与 rebuild\_down。

\begin{table}[htbp]
  \centering
  \caption{实验分组与主要验证内容}
  \label{tab:exp-plan}
  \small
  \begin{tabular}{p{1.6cm}p{3.6cm}p{5.8cm}}
    \toprule
    分组 & 负载方式 & 主要验证内容 \\
    \midrule
    规模 \(n\) & 均匀分布、混合增删 & 操作延迟、踢出、rebuild \\
    Facility \(K\) & 均匀分布、混合增删 & 操作延迟、踢出 \\
    k-kick & 均匀分布、高负载 & 操作延迟、踢出 \\
    Cubby & Insert/Delete/混合比例 & rebuild\_down、rebuild\_up、重构耗时 \\
    \bottomrule
  \end{tabular}
\end{table}

\subsection{评价指标}

为量化评估系统的基础性能，实验记录并分析以下指标：

\begin{itemize}
  \item 操作延迟：插入、命中查询和删除的平均延迟。该指标用于衡量哈希 系统基础操作的运行性能；
  \item k-kick：workload 内累计踢出次数，以及单次插入/删除搬移最大次数 insert\_moved\_max、delete\_moved\_max。该指标用于观察 k-kick 核心模块的具体表现，并判断模块的基础踢出功能是否生效；
  \item j-tiered Cubby：rebuild\_down（向下拆分）与 rebuild\_up（向上合并）触发次数。该指标用于验证多层级的 Cubby 的合并与拆分功能是否正常触发；
\end{itemize}

\section{规模 \texorpdfstring{$n$}{n} 实验}

固定 \(K=\log^3 n\)、\(k=4\)、\(c=2\)，令 \(n\) 从 \(10^3\) 增至 \(10^5\)。表~\ref{tab:exp-n-fixed} 给出各档测量结果。

\begin{table}[htbp]
  \centering
  \caption{规模 \(n\) 实验结果（\(K=\log^3 n\)，\(k=4\)，键均匀分布）}
  \label{tab:exp-n-fixed}
  \footnotesize
  \begin{tabular}{lrrrrrr}
    \toprule
    $n$ &
    ins ($\mu$s) & qry ($\mu$s) & del ($\mu$s) &
    总踢出 & rebuild\_down & rebuild\_up \\
    \midrule
    $10^3$  & 49.1 & 2.2  & 66.0 & 80     & 11    & 179   \\
    $10^4$  & 42.0 & 2.7  & 71.6 & 1157   & 77    & 930   \\
    $10^5$  & 60.8 & 5.1  & 2051.9 & 10485 & 11447 & 21502 \\
    \bottomrule
  \end{tabular}
\end{table}

\noindent 表中 rebuild\_down 对应向下拆分的次数，rebuild\_up 对应向上合并的次数。

\begin{figure}[htbp]
  \centering
  \begin{tikzpicture}
    \begin{axis}[
      width=0.88\linewidth,
      height=6.2cm,
      xmode=log,
      log basis x={10},
      xtick={1000,10000,100000},
      xticklabels={$10^3$,$10^4$,$10^5$},
      xmin=800, xmax=120000,
      xlabel={规模 $n$},
      ylabel={平均延迟 ($\mu$s)},
      legend style={font=\small},
      legend pos=north west,
      grid=major,
      grid style={dashed,gray!30},
    ]
    \addplot[mark=*, blue, thick] coordinates {(1000,49.1) (10000,42.0) (100000,60.8)};
    \addlegendentry{插入}
    \addplot[mark=square*, red, thick] coordinates {(1000,2.2) (10000,2.7) (100000,5.1)};
    \addlegendentry{查询}
    \addplot[mark=triangle*, green!60!black, thick] coordinates {(1000,66.0) (10000,71.6) (100000,2051.9)};
    \addlegendentry{删除}
    \end{axis}
  \end{tikzpicture}
  \caption{规模 $n$ 实验中三种操作平均延迟（$K=\log^3 n$，$k=4$，键均匀分布）}
  \label{fig:exp-n-latency}
\end{figure}

随 \(n\) 增大，命中查询延迟由 2.2\,$\mu$s 升至 5.1\,$\mu$s，相对三个数量级跨度增幅平缓，与 Local Query Router 的常数步定位一致\cite{bender2022otsh,morrison1968}。插入延迟在 42.0--60.8\,$\mu$s 之间，\(n=10^3\) 略高（49.1\,$\mu$s），与 Facility 数较少、单桶局部竞争有关。删除延迟由 66.0\,$\mu$s 增至 2051.9\,$\mu$s，增幅远高于插入：\(n=10^5\) 时 rebuild 合计约 3.3 万次，删除路径上的 tail 回填、tier 拆分及 MiniArray 同步均叠加在热路径上（见图~\ref{fig:exp-n-latency}）。累计踢出由 80 增至 \(1.0\times 10^4\)，rebuild\_down 由 11 增至 11447 次、rebuild\_up 由 179 增至 21502 次；大规模下 j-tiered Cubby 维护最为频繁\cite{bender2024iceberg,pandey2023iceberght}。

\section{Facility 规模 \texorpdfstring{$K$}{K} 实验}

固定 \(k=4\)，在 \(n\in\{10^3,10^4,10^5\}\) 三档下各取多组 \(K\) 比较三操作延迟与踢出。表~\ref{tab:exp-K-1e3}--\ref{tab:exp-K-1e5} 汇总实测结果（键均匀分布）。

\begin{table}[htbp]
  \centering
  \caption{Facility 规模 \(K\) 实验（\(n=10^3\)，\(k=4\)）}
  \label{tab:exp-K-1e3}
  \small
  \begin{tabular}{lrrrr}
    \toprule
    $K$ & ins ($\mu$s) & qry ($\mu$s) & del ($\mu$s) & 总踢出 \\
    \midrule
    16   & 22.6 & 2.1 & 52.8 & 112 \\
    64   & 49.9 & 2.0 & 67.7 & 60  \\
    256  & 51.9 & 1.9 & 66.0 & 58  \\
    \bottomrule
  \end{tabular}
\end{table}

\begin{table}[htbp]
  \centering
  \caption{Facility 规模 \(K\) 实验（\(n=10^4\)，\(k=4\)）}
  \label{tab:exp-K-fixed}
  \small
  \begin{tabular}{lrrrr}
    \toprule
    $K$ & ins ($\mu$s) & qry ($\mu$s) & del ($\mu$s) & 总踢出 \\
    \midrule
    64    & 56.9 & 2.8 & 75.9 & 813 \\
    256   & 54.1 & 2.7 & 70.9 & 916 \\
    1024  & 52.9 & 2.6 & 72.0 & 963 \\
    \bottomrule
  \end{tabular}
\end{table}

\begin{table}[htbp]
  \centering
  \caption{Facility 规模 \(K\) 实验（\(n=10^5\)，\(k=4\)）}
  \label{tab:exp-K-1e5}
  \small
  \begin{tabular}{lrrrr}
    \toprule
    $K$ & ins ($\mu$s) & qry ($\mu$s) & del ($\mu$s) & 总踢出 \\
    \midrule
    256   & 98.9  & 4.2 & 504.9  & 6889 \\
    1024  & 185.2 & 3.4 & 518.3  & 6883 \\
    4096  & 64.9  & 3.1 & 127.1  & 8756 \\
    \bottomrule
  \end{tabular}
\end{table}

\(n=10^3\) 时 \(K=16\) 插入最快（22.6\,$\mu$s），但踢出最多（112 次）；\(K\ge 64\) 时踢出降至约 60 次，增删延迟略升，体现小 \(n\) 下 Facility 过细反而增加路由与维护开销。\(n=10^4\) 三档 \(K\) 下查询均在 2.6--2.8\,$\mu$s，踢出随 \(K\) 增大变化不大（813--963），说明中等规模下 \(K\) 对均匀负载的边际影响有限。\(n=10^5\) 时 \(K=4096\) 删除延迟（127.1\,$\mu$s）明显低于 \(K=256/1024\)（504--518\,$\mu$s），与大规模 rebuild 主要落在较小 \(K\) 档、大 Facility 缓解局部冲突一致；\(K=1024\) 插入异常偏高（185.2\,$\mu$s）可能与 Facility 数与 Cubby 布局的组合有关，部署时宜结合目标 \(n\) 与负载率选取 \(K\)\cite{bender2022otsh,bender2024sosa}。

\section{k-kick 深度 \texorpdfstring{$k$}{k} 实验}

对 \(k\in\{3,4,5\}\) 在各规模下执行相同高负载混合增删 workload：\(n=10^3\) 取 \(K=64\)，\(n=10^4\) 取 \(K=256\)，\(n=10^5\) 取 \(K=4096\)。表~\ref{tab:exp-k-1e3}--\ref{tab:exp-k-1e5} 给出完整结果；图~\ref{fig:exp-k-bpk} 汇总三档规模下 bits/key 随 \(k\) 的变化。

\begin{table}[htbp]
  \centering
  \caption{k-kick 深度 \(k\) 实验（\(n=10^3\)，\(K=64\)，高负载）}
  \label{tab:exp-k-1e3}
  \small
  \begin{tabular}{lrrrr}
    \toprule
    $k$ & ins ($\mu$s) & qry ($\mu$s) & del ($\mu$s) & 总踢出 \\
    \midrule
    3 & 71.2 & 2.1 & 85.8 & 61  \\
    4 & 73.0 & 2.1 & 88.7 & 104 \\
    5 & 77.5 & 2.3 & 97.7 & 194 \\
    \bottomrule
  \end{tabular}
\end{table}

\begin{table}[htbp]
  \centering
  \caption{k-kick 深度 \(k\) 实验（\(n=10^4\)，\(K=256\)，高负载）}
  \label{tab:exp-k-fixed}
  \small
  \begin{tabular}{lrrrr}
    \toprule
    $k$ & ins ($\mu$s) & qry ($\mu$s) & del ($\mu$s) & 总踢出 \\
    \midrule
    3 & 54.5 & 2.7 & 72.6 & 265  \\
    4 & 55.3 & 2.6 & 72.6 & 900  \\
    5 & 56.4 & 2.5 & 70.2 & 1587 \\
    \bottomrule
  \end{tabular}
\end{table}

\begin{table}[htbp]
  \centering
  \caption{k-kick 深度 \(k\) 实验（\(n=10^5\)，\(K=4096\)，高负载）}
  \label{tab:exp-k-1e5}
  \small
  \begin{tabular}{lrrrr}
    \toprule
    $k$ & ins ($\mu$s) & qry ($\mu$s) & del ($\mu$s) & 总踢出 \\
    \midrule
    3 & 60.3 & 3.1 & 2229.4 & 2772  \\
    4 & 62.6 & 3.2 & 2003.9 & 8707  \\
    5 & 69.1 & 3.3 & 1734.3 & 14986 \\
    \bottomrule
  \end{tabular}
\end{table}


\begin{figure}[htbp]
  \centering
  \begin{tikzpicture}
    \begin{axis}[
      width=0.88\linewidth,
      height=6.2cm,
      xmin=2.5, xmax=5.5,
      xtick={3,4,5},
      xlabel={k-kick 深度 $k$},
      ylabel={bits/key},
      legend style={font=\small},
      legend pos=north east,
      grid=major,
      grid style={dashed,gray!30},
      ymin=184, ymax=200,
    ]
    \addplot[mark=*, blue, thick] coordinates {(3,191.19) (4,189.67) (5,186.53)};
    \addlegendentry{$n=10^3$}
    \addplot[mark=square*, red, thick] coordinates {(3,192.20) (4,189.10) (5,186.12)};
    \addlegendentry{$n=10^4$}
    \addplot[mark=triangle*, green!60!black, thick] coordinates {(3,198.39) (4,195.48) (5,192.65)};
    \addlegendentry{$n=10^5$}
    \end{axis}
  \end{tikzpicture}
  \caption{k-kick 深度 $k$ 实验中 bits/key}
  \label{fig:exp-k-bpk}
\end{figure}

\noindent insert\_moved\_max 与 delete\_moved\_max 各档均不超过 \(k\)。

三档规模下累计踢出均随 \(k\) 增大而上升：\(n=10^3\) 由 61 增至 194，\(n=10^4\) 由 265 增至 1587，\(n=10^5\) 由 2772 增至 14986\cite{bender2022otsh,pagh2004cuckoo}。查询延迟在各规模均不超过 3.3\,$\mu$s，kick 深度主要作用于增删路径。\(n=10^5\) 删除延迟达毫秒级，且 \(k=5\) 时（1734.3\,$\mu$s）低于 \(k=3/4\)，与大规模 rebuild 叠加、较深 kick 链在特定负载下分摊搬移的实测表现一致；中小规模下增删延迟随 \(k\) 变化相对平缓\cite{bender2022otsh,bender2024sosa}。图~\ref{fig:exp-k-bpk},表明，随着 k-kick 树深度的增加，每键占用的内存呈现下降趋势，这与论文核心观点相符。

\section{多层级 Cubby 动态重构实验}

固定 \(k=4\)，在 \(n\in\{10^3,10^4,10^5\}\) 下分别取 \(K=64\)、\(256\)、\(4096\)，按表~\ref{tab:exp-tier-1e3}--\ref{tab:exp-tier-1e5} 四种 workload 执行操作序列，统计 rebuild 次数与单次重构平均耗时。

\begin{table}[htbp]
  \centering
  \caption{多层级 Cubby 动态重构（\(n=10^3\)，\(K=64\)）}
  \label{tab:exp-tier-1e3}
  \small
  \begin{tabular}{lrrr}
    \toprule
    Workload      & rebuild\_down & rebuild\_up & 平均重构耗时 (ms) \\
    \midrule
    Insert-only   & 1   & 31  & 3.72 \\
    Delete-only   & 28  & 0   & 6.24 \\
    80/20 混合    & 0   & 24  & 3.71 \\
    50/50 混合    & 6   & 25  & 3.82 \\
    \bottomrule
  \end{tabular}
\end{table}

\begin{table}[htbp]
  \centering
  \caption{多层级 Cubby 动态重构（\(n=10^4\)，\(K=256\)）}
  \label{tab:exp-tier-fixed}
  \small
  \begin{tabular}{lrrr}
    \toprule
    Workload      & rebuild\_down & rebuild\_up & 平均重构耗时 (ms) \\
    \midrule
    Insert-only   & 0   & 384 & 3.47 \\
    Delete-only   & 384 & 2   & 6.10 \\
    80/20 混合    & 8   & 355 & 3.32 \\
    50/50 混合    & 68  & 377 & 3.35 \\
    \bottomrule
  \end{tabular}
\end{table}

\begin{table}[htbp]
  \centering
  \caption{多层级 Cubby 动态重构（\(n=10^5\)，\(K=4096\)）}
  \label{tab:exp-tier-1e5}
  \small
  \begin{tabular}{lrrr}
    \toprule
    Workload      & rebuild\_down & rebuild\_up & 平均重构耗时 (ms) \\
    \midrule
    Insert-only   & 0    & 3591 & 5.64 \\
    Delete-only   & 4211 & 621  & 9.76 \\
    80/20 混合    & 177  & 3588 & 5.52 \\
    50/50 混合    & 375  & 3504 & 5.13 \\
    \bottomrule
  \end{tabular}
\end{table}

三档规模下规律一致：纯插入以 rebuild\_up 为主（\(n=10^3\) 为 31 次，\(n=10^5\) 为 3591 次），纯删除以 rebuild\_down 为主（28--4211 次），对应 Cubby 随元素增减而向上合并或向下拆分\cite{bender2024iceberg,pandey2023iceberght}。80/20 混合负载在各规模均以 rebuild\_up 为主；50/50 混合时两方向均显著（如 \(n=10^4\) 为 68/377，\(n=10^5\) 为 375/3504），删除比例升高使 rebuild\_down 占比上升。平均重构耗时由 3.3--6.2\,ms（\(n\le 10^4\)）增至 5.1--9.8\,ms（\(n=10^5\)），仍远低于大规模混合 workload 中单次删除的均摊延迟。

\section{本章小结}

本章在 \(n\)、Facility 规模 \(K\)、k-kick 深度 \(k\) 及 Cubby 动态重构四组设定下，于 \(n\in\{10^3,10^4,10^5\}\) 完成测量。表~\ref{tab:exp-n-fixed} 与图~\ref{fig:exp-n-latency} 显示查询延迟由 2.2\,$\mu$s 升至 5.1\,$\mu$s，\(n=10^5\) 删除延迟达 2051.9\,$\mu$s；\(K\)、\(k\) 与 Cubby 三组实验在表~\ref{tab:exp-K-1e3}--\ref{tab:exp-tier-1e5} 中给出三档规模对照，insert\_moved\_max 各档均不超过 \(k\)，纯插入/纯删除 workload 分别主导 rebuild\_up 与 rebuild\_down。上述结果与 Local Query Router 常数步查询设计及 k-kick 交换上界在实测中可对照。